% =====================================================================
%  Cover, title page and back cover of the pocket dictionary.
%
%  The model is the mock-up supplied by the publisher of this edition
%  (Dicionario.pages, A4): front cover blue with black lettering and a
%  white star, title page white with black lettering and a blue star,
%  back cover plain blue. Everything that follows -- type sizes, line
%  heights, diameter of the emblem -- is SURVEYED on that model to the
%  300th of an inch, then carried over to the pocket page. The measured
%  values are given alongside, so that the calculation can be redone.
% =====================================================================

\usepackage{tikz}
\usepackage{eso-pic}

% The blue is AZURE, #007FFF -- the reference blue, exactly halfway
% between blue and cyan on the sRGB wheel (hue 210 degrees, saturation
% and value full). It is the blue of the two covers, of the disc and of
% the ID lettering of the title page: one single colour for the whole
% book.
%
% It does not come from a survey of the mock-up: that one is written in
% the Display P3 space, where its blue is (0.2549 ; 0.5728 ; 0.9688),
% that is #0094FF once converted to sRGB by the profile enclosed in the
% file. Azure is close to it -- twenty-one points of green lower -- and
% it is a NAMED colour, not the residue of a conversion.
\definecolor{idoblua}{RGB}{0,127,255}

% Futura is not free; Jost* is its revival under an open licence.
% Loaded BY PATH, not by name: the document sets identically on a
% machine where the font is not installed. See fonts/README.md.
\newfontfamily\kovrilfonto{Jost}[Path=fonts/, Extension=.ttf,
  UprightFont=*-Medium, BoldFont=*-Bold, ItalicFont=*-MediumItalic]

% --- Type sizes -------------------------------------------------------
% The model is set in Futura on an A4 page (595.28 pt): 62 pt for
% DICIONARIO, 38 pt for the subtitle, 204 pt for IDO, 29 pt for the
% signature. Two corrections carry them over to the pocket page:
%
%   1. the ratio of the page widths, 110 mm / 210 mm;
%   2. the ratio of the capital heights, Futura 0.754 em against
%      Jost 0.700 em, that is 1.0770 -- without which the lettering, on
%      a cover that is almost nothing but capitals, would look smaller
%      than the model at equal size.
%
% The two compose into a single factor, applied to \paperwidth.
\newcommand{\korpDICIONARIO}{0.11217\paperwidth}   % 62 pt sur A4
\newcommand{\korpSUBTITOLO} {0.06875\paperwidth}   % 38 pt
\newcommand{\korpIDO}       {0.36908\paperwidth}   % 204 pt
\newcommand{\korpNOMO}      {0.05247\paperwidth}   % 29 pt

% --- Line heights -----------------------------------------------------
% Position of the TOP of the ink of each line, counted from the top edge
% of the sheet. Three decisions fix them.
%
% 1. THE SCALE. The vertical distances follow the ratio of the page
%    WIDTHS, 110/210 = 0.5238, and not that of the type sizes. The two
%    differ: the size was enlarged by 1.0770 to equalise the capital
%    height of Jost* and that of Futura, and it is the capital, not the
%    em, that the eye measures. A vertical distance scaled to the
%    capitals is therefore scaled to the widths -- the two factors
%    cancel exactly, and the diameter of the emblem confirms it:
%    204 x 0.5641 x 0.700 = 153.8 x 0.5238.
%
% 2. THE LEADING OF THE SUBTITLE. The model carries 50.1 pt between the
%    first and the second line, 56.2 between the second and the third,
%    where all three are of the same size. That is a paragraph setting
%    left uneven, not an intention: we take the mean, 53.15 pt, and
%    « di la » comes down 3.05 pt on the A4 template, that is 1.6 pt on
%    the pocket page.
%
% 3. THE PLACING OF THE BLOCK. The two margins that remain are divided
%    in GOLDEN SECTION, the small one at the top: 81.75 pt against
%    132.27, whose ratio is 1/phi = 0.6180. The model divided them 87.8
%    against 188.6, that is 0.466 -- a proportion that suited an A4 page
%    but which, carried over to a page taller than it is wide, pushed
%    the composition too high.
%
%    The check comes from the emblem, which is what the eye weighs: its
%    golden section centres it at 0.5297 of the page height, where the
%    model places it at 0.5323 -- two thousandths apart. The old placing
%    put it at 0.5034, that is, at the dead middle of the page.
\newcommand{\posTITOLO}  {0.160218\paperheight}
\newcommand{\posRADIKI}  {0.235673\paperheight}
\newcommand{\posDILA}    {0.290237\paperheight}
\newcommand{\posLINGUO}  {0.344801\paperheight}
\newcommand{\posIDO}     {0.445613\paperheight}
\newcommand{\posNOMO}    {0.675162\paperheight}
\newcommand{\posAKADEMIO}{0.715404\paperheight}

% --- The emblem -------------------------------------------------------
% A disc and a six-pointed star. The star is REGULAR, and is built:
%
%   1. an equilateral triangle at the centre, one point down -- its
%      three vertices are the three short points;
%   2. on each of its sides, a long point raised outwards, whose tip
%      touches the circle;
%   3. the long point at the top has the same height as the triangle;
%   4. each long point has for its base ONE THIRD of the side that
%      carries it.
%
% Rule 4 covers another, more telling, and probably the rule of the
% original symbol: the edge of a long point, EXTENDED inwards, passes
% through the middle of one of the other two sides of the triangle. The
% two rules give the same figure, and it is again t = 1/2 that reconciles
% them -- the construction by the midpoints gives the base (2-t)/(4+t) of
% the side, that is exactly a third when t is one half.
%
% Condition 3 alone fixes everything. Taking the radius of the disc for
% the unit and t for the circumscribed radius of the triangle, the height
% of the triangle is 3t/2 and that of a point is 1 minus the inscribed
% radius, that is 1 - t/2; equating them gives t = 1/2. The triangle
% therefore fits in the half-disc, its side is V3/2, its inscribed radius
% 1/4, and a third of its side V3/6 -- whence all the numbers below,
% which are multiples of V3/12 and of 1/4.
%
% The figure is invariant under a rotation of a third of a turn: the
% three long points are on the circle, the three short ones at half the
% radius.
%
% Set against the mock-up, whose star is drawn by hand, the constructed
% figure departs from it by 2.4 % of the radius at most -- the star of
% the mock-up sits 1.4 % too low in its disc, and its long points have at
% the base 0.30 times the side and not a third.
%
% Coordinates in units of RADIUS, origin at the centre of the disc, y
% upwards. The baseline of the figure passes through the centre of the
% disc.
\newlength{\stelrayo}
\newcommand{\idoemblemo}[3]{% #1 diametre, #2 couleur du disque, #3 de l'etoile
  \setlength{\stelrayo}{0.5\dimexpr#1\relax}%
  \begin{tikzpicture}[x=\stelrayo, y=\stelrayo, baseline={(0,0)}]
    \fill[#2] (0,0) circle[radius=1];
    \fill[#3]
         (            0 ,  1   )   % longue pointe du haut
      -- ({-sqrt(3)/12 },  0.25)   % base de la pointe du haut
      -- ({-sqrt(3)/4  },  0.25)   % sommet gauche du triangle
      -- ({-sqrt(3)/6  },  0   )   % base de la pointe bas-gauche
      -- ({-sqrt(3)/2  }, -0.5 )   % longue pointe bas-gauche
      -- ({-sqrt(3)/12 }, -0.25)   % base de la pointe bas-gauche
      -- (            0 , -0.5 )   % sommet bas du triangle
      -- ({ sqrt(3)/12 }, -0.25)
      -- ({ sqrt(3)/2  }, -0.5 )   % longue pointe bas-droite
      -- ({ sqrt(3)/6  },  0   )
      -- ({ sqrt(3)/4  },  0.25)   % sommet droit du triangle
      -- ({ sqrt(3)/12 },  0.25)
      -- cycle;
  \end{tikzpicture}}

% The word IDO: the first two letters are set, the third is the emblem.
% On the model the disc is 1.0651 times the capital height and centres on
% it, and it follows the D at 1.4 pt.
%
% The letters take the colour of the DISC and not that of the text: on
% the cover they are white when the title is black, on the title page
% they are blue. The word makes one block with the emblem, it does not
% make one block with the lettering.
\newlength{\kapalto}
% The diameter, measured once and for all at the opening of the
% document: the two covers and the title page carry the SAME emblem, and
% it must not depend on where one lays it.
\newlength{\emblemdiam}
\AtBeginDocument{{\kovrilfonto\bfseries\fontsize{\korpIDO}{\korpIDO}\selectfont
  \settoheight{\kapalto}{ID}\global\emblemdiam=1.0651\kapalto}}
\newcommand{\idogrupo}[2]{% #1 couleur du disque, #2 couleur de l'etoile
  {\kovrilfonto\bfseries\fontsize{\korpIDO}{\korpIDO}\selectfont
   \settoheight{\kapalto}{ID}%
   \textcolor{#1}{ID}\kern0.0045\paperwidth
   \raisebox{0.5\kapalto}{\idoemblemo{\emblemdiam}{#1}{#2}}}}

% --- Placing the lines ------------------------------------------------
% \supre{height}{content}: the content centred on the width of the
% sheet, the TOP of its ink at the stated height. \raisebox{-\height}
% lowers the box by its own height, which brings its top onto the
% baseline of the \put.
\newcommand{\supre}[2]{%
  \put(0,-#1){\makebox[\paperwidth][c]{\raisebox{-\height}{#2}}}}

% The lettering, identical on both pages; only the colours change.
\newcommand{\kovrilteksto}[3]{% #1 encre, #2 disque, #3 etoile
  \AddToShipoutPictureFG*{\AtPageUpperLeft{%
    \setlength{\unitlength}{1pt}\color{#1}\kovrilfonto
    \supre{\posTITOLO}{\bfseries\fontsize{\korpDICIONARIO}{\korpDICIONARIO}
      \selectfont DICIONARIO}%
    \supre{\posRADIKI}{\bfseries\fontsize{\korpSUBTITOLO}{\korpSUBTITOLO}
      \selectfont de la 10.000 RADIKI}%
    \supre{\posDILA}{\bfseries\fontsize{\korpSUBTITOLO}{\korpSUBTITOLO}
      \selectfont di la}%
    \supre{\posLINGUO}{\bfseries\fontsize{\korpSUBTITOLO}{\korpSUBTITOLO}
      \selectfont linguo universala}%
    \supre{\posIDO}{\idogrupo{#2}{#3}}%
    \supre{\posNOMO}{\fontsize{\korpNOMO}{\korpNOMO}\selectfont
      Marcelo Persiko \textit{(Marcel Pesch)}}%
    \supre{\posAKADEMIO}{\fontsize{\korpNOMO}{\korpNOMO}\selectfont
      ex-membro di la Akademio}%
  }}}

\newcommand{\fondo}[1]{%
  \AddToShipoutPictureBG*{\AtPageLowerLeft{%
    \color{#1}\rule{\paperwidth}{\paperheight}}}}

% --- The three pages --------------------------------------------------
% The front cover: blue ground, black lettering, white disc, blue star.
\newcommand{\platrekto}{%
  \fondo{idoblua}\kovrilteksto{black}{white}{idoblua}%
  \thispagestyle{empty}\null\clearpage}

% The title page: white ground, black lettering, blue disc, white star.
\newcommand{\titolpagino}{%
  \kovrilteksto{black}{idoblua}{white}%
  \thispagestyle{empty}\null\clearpage}

% The back cover: the emblem alone on the blue. The same diameter as on
% the front -- it is the same mark, and repeating it at the same size
% makes the two covers rhyme. The disc is white and the star blue, as on
% the front: on a blue ground it is the disc that must carry the light.
%
% The height is not the geometric middle: a sign on its own looks LOW
% there, and the optical centre is wanted. The one the eye asks for here
% happens to be the MIRROR of the centre of the emblem of the front cover
% about the horizontal axis of the page -- 0.5297 on the front, 1 minus
% 0.5297 on the back. The two covers therefore hold the mark at
% symmetrical heights: a little below the middle when a composition
% accompanies it, a little above when it is alone.
\newcommand{\posVERSO}{0.470298\paperheight}
\newlength{\posverso}
\newcommand{\platverso}{%
  \fondo{idoblua}%
  \AddToShipoutPictureFG*{\AtPageUpperLeft{%
    \setlength{\unitlength}{1pt}%
    \setlength{\posverso}{\posVERSO}%
    \addtolength{\posverso}{-0.5\emblemdiam}%
    \put(0,-\posverso){\makebox[\paperwidth][c]{\raisebox{-\height}{%
      \idoemblemo{\emblemdiam}{white}{idoblua}}}}}}%
  \thispagestyle{empty}\null\clearpage}
